% from ega0-20-fr.tex
\oldpage[0\textsubscript{IV-1}]{142}
de $C$-modules $\mathrm{K}_\bullet(C/A/\Lambda)$ et
$\mathrm{K}_\bullet(C/B/\Lambda)$ (20.6.5), et en outre les complexes de
$C$-modules suivants. On posera tout d'abord
\begin{equation}
\tag{20.6.15.1}
\label{0.20.6.15.1-fr}
  \mathrm{K}^C_\bullet(B/A/\Lambda)=
  \mathrm{K}_\bullet(B/A/\Lambda)\otimes_B C.
\end{equation}
D'autre part, nous désignerons par $\mathrm{T}_\bullet(C/B/\Lambda)$ le
complexe de $C$-modules dont les termes sont nuls sauf en degrés 0 et 1,
et où
\begin{equation}
\tag{20.6.15.2}
\label{0.20.6.15.2-fr}
  \mathrm{T}_0(C/B/\Lambda)=\mathrm{T}_1(C/B/\Lambda)=
  \Omega^1_{B/\Lambda}\otimes_B C
\end{equation}
l'opérateur différentiel étant l'identité, de sorte que ce complexe est
homotope à 0; posons enfin
\begin{equation}
\tag{20.6.15.3}
\label{0.20.6.15.3-fr}
  \mathrm{K}'_\bullet(C/A/\Lambda)=
  \mathrm{K}_\bullet(C/A/\Lambda)\oplus
  \mathrm{T}_\bullet(C/B/\Lambda).
\end{equation}
En vertu du caractère trivial de $\mathrm{T}_\bullet$, il est clair que
l'on a
\begin{equation}
\tag{20.6.15.4}
\label{0.20.6.15.4-fr}
  \mathrm{H}^i(\mathrm{K}'_\bullet,L)=
  \mathrm{H}^i(\mathrm{K}_\bullet,L)
  \quad\text{et}\quad
  \mathrm{H}_i(\mathrm{K}'_\bullet,L)=
  \mathrm{H}_i(\mathrm{K}_\bullet,L)
\end{equation}
pour tout $C$-module $L$ et tout $i$.
\end{env}
\begin{env}[20.6.16]
\label{0.20.6.16-fr}
Définissons maintenant une suite exacte de complexes, \emph{scindée} en chaque
degré
\begin{equation}
\tag{20.6.16.1}
\label{0.20.6.16.1-fr}
  0\to\mathrm{K}^C_\bullet(B/A/\Lambda)
  \xrightarrow{j}\mathrm{K}'_\bullet(C/A/\Lambda)
  \xrightarrow{p}\mathrm{K}_\bullet(C/B/\Lambda)\to0
\end{equation}
de la façon suivante : désignons pour un moment par
\[
  f:\Omega^1_{A/\Lambda}\otimes_A C\to
    \Omega^1_{B/\Lambda}\otimes_B C,
  \qquad
  g:\Omega^1_{B/\Lambda}\otimes_B C\to\Omega^1_{C/\Lambda}
\]
les homomorphismes canoniques $u_{B/A/\Lambda}\otimes1_C$ et
$v_{C/B/\Lambda}$ respectivement, dont le composé est
$g\circ f=(v\circ u)_{C/A/\Lambda}$ (cf.\ \sref{0.20.6.14.4-fr}).
On prendra $j_1(x)=(x,f(x))$, $p_1(y,z)=z-f(y)$,
$j_0(x)=(g(x),x)$, $p_0(y,z)=g(z)-y$ de sorte que
$\operatorname{Im}(j_1)=\operatorname{Ker}(p_1)$ est le graphe de $f$,
supplémentaire de $\{0\}\oplus\mathrm{T}_1$, et
$\operatorname{Im}(j_0)=\operatorname{Ker}(p_0)$ est le graphe de $g$,
supplémentaire de $\mathrm{K}_1(C/A/\Lambda)\oplus\{0\}$ ; la vérification
de la commutativité du diagramme
\[
\xymatrix{
  0 \ar[r] & \mathrm{K}^C_1(B/A/\Lambda) \ar[r]^{j_1} \ar[d] &
  \mathrm{K}'_1(C/A/\Lambda) \ar[r]^{p_1} \ar[d] &
  \mathrm{K}_1(C/B/\Lambda) \ar[r] \ar[d] & 0 \\
  0 \ar[r] & \mathrm{K}^C_0(B/A/\Lambda) \ar[r]_{j_0} &
  \mathrm{K}'_0(C/A/\Lambda) \ar[r]_{p_0} &
  \mathrm{K}_0(C/B/\Lambda) \ar[r] & 0
}
\]
où les flèches verticales sont les opérateurs différentiels, est immédiate.
\end{env}

\begin{theorem}[20.6.17]
\label{0.20.6.17-fr}
On a une suite exacte de $C$-modules
\begin{equation}
\tag{20.6.17.1}
\label{0.20.6.17.1-fr}
  0\to\Upsilon^C_{B/A/\Lambda}\xrightarrow{v'}
  \Upsilon_{C/A/\Lambda}\xrightarrow{u'}
  \Upsilon_{C/B/\Lambda}\xrightarrow{\partial}
  \Omega^1_{B/A}\otimes_B C\xrightarrow{v_{C/B/A}}
  \Omega^1_{C/A}\xrightarrow{u_{C/B/A}}\Omega^1_{C/B}\to0
\end{equation}
